% Self-contained article sections.  Standard packages: amsmath, amssymb,
% amsthm, mathtools.  Standard theorem environments are used.
% Macros expected from parent: \Li, \dd, \ii, \Ree, \C, \N.
% No additional custom macros are required.  Labels use the prefix sfj:.

\section{Exact coefficients for every spectral face}
\label{sfj:sec:faces}

The newest incoming spectral-face report gives an all-order common-ray
formula and asks, at the end of its spectral section, for complete
collision coefficients and independently variable spectral orders
\cite{RepoBernoulli}.  We answer both questions.  A further
specialization evaluates every face with just one displaced mass by
an ordinary series of \emph{depth-two} polylogarithms.  The classical
Dirichlet integral, Gamma product, and polylogarithm differentiation
identities are ingredients, not novelty claims
\cite{Carlson,DLMF}.

\subsection{The inherited face normalization}

Let $k\ge2$, let $b_1,\ldots,b_k$ be real, and let
\[
 \Sigma_k=\{t_i\ge0:\ t_1+\cdots+t_k=1\}.
\]
The measure is $dt_1\cdots dt_{k-1}$, with the last coordinate
eliminated.  For a nonempty subset $T\subseteq\{1,\ldots,k\}$ put
\[
 t_T=\sum_{i\in T}t_i,\qquad
 b_T=\frac{\sum_{i\in T}t_i b_i}{t_T}.
\]
The perspective $t_TF(b_T)$ is filled by zero when $t_T=0$.  Define
\begin{align}
 \mathcal H_k(F;t)&=
 \sum_{\varnothing\ne T\subseteq\{1,\ldots,k\}}
       (-1)^{k-|T|}t_TF(b_T),\label{sfj:eq:face}\\
 I_k(F;b_1,\ldots,b_k)&=
 \int_{\Sigma_k}\mathcal H_k(F;t)\frac{dt}{t_1\cdots t_k}.
 \label{sfj:eq:I}
\end{align}
The notation $I_J$ denotes the same construction with the labels in a
subset $J$.  This is exactly the predecessor's face normalization; in
particular no Euclidean surface-area factor is inserted.

For $F\in C^1([\min b_i,\max b_i])$, the integral is ordinary and
absolutely convergent.  Indeed the perspective is Lipschitz on its
nonnegative coordinate cone, because
\[
 \partial_{t_i}\{t_TF(b_T)\}
      =F(b_T)+F'(b_T)(b_i-b_T)\quad(i\in T).
\]
Terms cancel in pairs when any coordinate is zero.  Consequently
$|\mathcal H_k(F;t)|\le C\min_i t_i$.  On a region where one
coordinate is maximal, it is at least $1/k$ and
$\min_i t_i\le(\prod_{i\ne i_{\max}}t_i)^{1/(k-1)}$.
This also proves that
\begin{equation}\label{sfj:eq:normal-face-domain}
 \int_{\Sigma_k}\mathcal H_k(F;t)\prod_i t_i^{\alpha_i-1}\,dt
 \quad\hbox{is holomorphic if}\quad
 \Ree\alpha_i>-\frac1{k-1}\ \hbox{for every }i,
\end{equation}
with locally normally convergent logarithmic derivatives.  These
elementary estimates justify every face integral used below.  Constants
and affine functions are annihilated by $\mathcal H_k$.

\subsection{A finite harmonic formula at every polynomial degree}

For $m\ge2$ and $\ell\ge0$ define
\begin{equation}\label{sfj:eq:complete-harmonic}
 \mathfrak h_\ell(m)
  =[v^\ell]\prod_{j=1}^{m-1}(1-v/j)^{-1}.
\end{equation}
Thus $\mathfrak h_\ell(m)$ is the complete homogeneous polynomial
of degree $\ell$ in $1,1/2,\ldots,1/(m-1)$.  In the usual generalized
harmonic notation, its first values are
\[
 1,\qquad H_{m-1},\qquad
 \frac{H_{m-1}^2+H_{m-1}^{(2)}}2,\qquad
 \frac{H_{m-1}^3+3H_{m-1}H_{m-1}^{(2)}+2H_{m-1}^{(3)}}6.
\]
It is also the weakly ordered harmonic sum
$H_{m-1}^{\star}(\{1\}^{\ell})$.

\begin{theorem}[Every monomial face integral]
\label{sfj:thm:monomial}
Let $k\ge2$ and $m\ge2$.  For a vector
$\boldsymbol h=(h_1,\ldots,h_k)\in\N_0^k$ let
$d(\boldsymbol h)=|\{i:h_i>0\}|$.  Then
\begin{equation}\label{sfj:eq:monomial}
 \boxed{\displaystyle
 I_k(B\mapsto B^m;\boldsymbol b)
  =m\!\sum_{h_1+\cdots+h_k=m}
   (-1)^{k-d}(k-d)!\,\mathfrak h_{k-d}(m)
                  \prod_{h_i>0}\frac{b_i^{h_i}}{h_i}.}
\end{equation}
The right side is a finite symmetric polynomial with rational
coefficients.  The formula extends to complex $b_i$ as a polynomial
identity.
\end{theorem}

\begin{proof}
The perspective is homogeneous of degree one.  Radial integration
therefore gives
\[
 I_k(F;\boldsymbol b)=
 \int_{(0,\infty)^k} e^{-\sum_i t_i}
            \mathcal H_k(F;t)\frac{dt_1\cdots dt_k}{\prod_i t_i};
\]
the radial factor is $\int_0^\infty e^{-\rho}\,d\rho=1$.
For a monomial of degree $m\ge2$, use
\[
 t_T b_T^m=\frac{(\sum_{i\in T}t_i b_i)^m}{t_T^{m-1}},\qquad
 t_T^{1-m}=\frac1{(m-2)!}\int_0^\infty
                      v^{m-2}e^{-vt_T}\,dv.
\]
Fix a monomial vector $\boldsymbol h$ with support $U$ and $d=|U|$.
The alternating sum over $T\supseteq U$ is
\[
 e^{-v\sum_{i\in U}t_i}
             \prod_{j\notin U}(e^{-vt_j}-1).
\]
Integrating the active and inactive coordinates respectively gives
\[
 \frac{\Gamma(h_i)}{(1+v)^{h_i}},\qquad
 \int_0^\infty e^{-t}\frac{e^{-vt}-1}{t}\,dt=-\log(1+v).
\]
The last identity follows by differentiating in $v$ and fixing the
value at zero.  Every resulting term is absolutely integrable: its
remaining majorant is a constant times
$v^{m-2}(1+v)^{-m}\log^{k-d}(1+v)$.  Thus the use of Fubini is
legitimate after the finite inclusion--exclusion has been combined.
We obtain
\begin{align*}
 I_k(B^m;\boldsymbol b)
  =m(m-1)\sum_{|\boldsymbol h|=m}
     (-1)^{k-d}\prod_{h_i>0}\frac{b_i^{h_i}}{h_i}
       \int_0^\infty
        \frac{v^{m-2}\log^{k-d}(1+v)}{(1+v)^m}\,dv.
\end{align*}
On setting $x=(1+v)^{-1}$, the integral becomes
\[
 \int_0^1(1-x)^{m-2}(-\log x)^{\ell}\,dx
       =\frac{\ell!}{m-1}\mathfrak h_\ell(m),\qquad \ell=k-d.
\]
To verify the last equality, expand the elementary beta integral
\[
 \int_0^1x^u(1-x)^{m-2}\,dx
 =\frac1{m-1}\prod_{j=1}^{m-1}(1+u/j)^{-1}
\]
at $u=0$.  Substitution proves the stated formula.
\end{proof}

At $m=2$, Theorem~\ref{sfj:thm:monomial} reduces to the predecessor's
factorial variance identity
\begin{equation}\label{sfj:eq:quadratic}
 I_k(B^2;\boldsymbol b)
   =(-1)^{k-1}(k-2)!\sum_{i<j}(b_i-b_j)^2.
\end{equation}
The new formula supplies every degree beyond that quadratic term.  For
example,
\begin{equation}\label{sfj:eq:three-polynomials}
 I_3(B^2;0,1,2)=6,\qquad I_3(B^3;0,1,2)=18,\qquad
 I_3(B^4;0,1,2)=\frac{97}{2},\qquad
 I_3(B^5;0,1,2)=\frac{245}{2}.
\end{equation}

\subsection{Complete convergent collision expansions}

\begin{corollary}[All collision coefficients]
\label{sfj:cor:collision}
Suppose $F$ is holomorphic on $|z-B|<R$ and let
$M=\max_i|\beta_i|$.  For $|\epsilon|M<R$,
\begin{equation}\label{sfj:eq:collision}
 I_k(F;B+\epsilon\beta_1,\ldots,B+\epsilon\beta_k)
  =\sum_{m=2}^\infty
       \frac{F^{(m)}(B)}{m!}\epsilon^m
                     I_k(B^m;\boldsymbol\beta),
\end{equation}
normally on compact subsets of this disk.  Every coefficient is given
by the finite harmonic formula \eqref{sfj:eq:monomial}.
\end{corollary}
\begin{proof}
The constant and affine Taylor terms vanish.  On each smaller closed
disk, the Taylor series and its first derivative converge uniformly.
The boundary estimate preceding \eqref{sfj:eq:normal-face-domain}
is a continuous $C^1$ estimate and permits termwise integration.
Cauchy estimates on an intermediate disk give a geometric majorant,
including the harmless linear factor in $m$ from the derivative.
\end{proof}

The Gamma example is especially explicit.  For $a>0$ and $B\ge0$ put
\begin{align}
 G_a(B)&=\log\Gamma(a+\ii\sqrt B)
          +\log\Gamma(a-\ii\sqrt B)-\log(2\pi),
                    \label{sfj:eq:gammaG}\\
 R_m(a;B)&=\sum_{n=0}^\infty\bigl((n+a)^2+B\bigr)^{-m}.
                    \label{sfj:eq:resolvent}
\end{align}
The Gamma logarithms are analytic logarithms continued from positive
real arguments.  The paired function is analytic as a function of
$B$ on the disk of radius $a^2+B_0$ about any $B_0\ge0$.  The
Gamma product, or differentiation of its normally convergent logarithmic
difference, gives
\begin{equation}\label{sfj:eq:gamma-derivatives}
 G_a^{(m)}(B)=(-1)^m(m-1)!R_m(a;B),\qquad m\ge1.
\end{equation}
Consequently
\begin{equation}\label{sfj:eq:gamma-collision}
 \boxed{\displaystyle
 I_k(G_a;B+\epsilon\boldsymbol\beta)
   =\sum_{m=2}^\infty\frac{(-1)^m\epsilon^m}{m}
                  R_m(a;B)I_k(B^m;\boldsymbol\beta),
 \quad |\epsilon|M<a^2+B.}
\end{equation}
At $B=0$ the resolvent is simply $\zeta(2m,a)$.  For instance,
\begin{equation}\label{sfj:eq:gamma-three-example}
 I_3(G_a;0,\epsilon,2\epsilon)
   =3\zeta(4,a)\epsilon^2-6\zeta(6,a)\epsilon^3
       +\frac{97}{8}\zeta(8,a)\epsilon^4
       -\frac{49}{2}\zeta(10,a)\epsilon^5+\cdots,
 \qquad 2|\epsilon|<a^2.
\end{equation}
The positive integer resolvents themselves have a finite polygamma
description.  For $B>0$,
\begin{equation}\label{sfj:eq:resolvent-digamma}
 R_1(a;B)=
 \frac{\psi(a+\ii\sqrt B)-\psi(a-\ii\sqrt B)}{2\ii\sqrt B},
 \qquad
 R_m(a;B)=\frac{(-1)^{m-1}}{(m-1)!}
                 \partial_B^{m-1}R_1(a;B).
\end{equation}
Their apparent singularities at $B=0$ are removable.

\subsection{All logarithmic face moments have a rational compiler}

For independent parameters $\alpha_1,\ldots,\alpha_k$ define
\[
 A=\sum_i\alpha_i,\qquad
 Q_k(\boldsymbol\alpha)=
       \frac{\Gamma(1+A)}{\prod_i\Gamma(1+\alpha_i)},\qquad
 K_{k,m}(\boldsymbol\alpha;\boldsymbol b)
   =Q_k(\boldsymbol\alpha)
       \int_{\Sigma_k}\mathcal H_k(B^m;t)
                                      \prod_i t_i^{\alpha_i-1}\,dt.
\]
For an increment $v$, let
$\Delta_v f(y)=(f(y+v)-f(y))/v$, with $\Delta_0f=f'$.

\begin{theorem}[Rational normalized face moments]
\label{sfj:thm:rational}
Let $m\ge2$ and $f_m(y)=(1+y)_{m-1}^{-1}$.  For
$U=\{i:h_i>0\}$, write $A_U=\sum_{i\in U}\alpha_i$.  Then
\begin{equation}\label{sfj:eq:rational}
 \boxed{\displaystyle
 K_{k,m}=m!\sum_{|\boldsymbol h|=m}
  \left(\prod_{i=1}^k\frac{b_i^{h_i}}{h_i!}\right)
  \left(\prod_{i\in U}(1+\alpha_i)_{h_i-1}\right)
  \left(\prod_{j\notin U}\Delta_{\alpha_j}\right)f_m(A_U).}
\end{equation}
In particular this Gamma-normalized integral is a rational function of
all the spectral parameters, holomorphic in a neighborhood of zero.
Its Taylor coefficients give all its logarithmically weighted face
moments by finite rational algebra.  Removing $Q_k$ introduces only
the finite zeta coefficients in the Taylor series of $\log\Gamma$.
\end{theorem}
\begin{proof}
Repeat the orthant proof of Theorem~\ref{sfj:thm:monomial} with powers
$t_i^{\alpha_i-1}$.  The radial factor is $\Gamma(1+A)$, and the
normalization leaves division by $\prod_i\Gamma(1+\alpha_i)$.
For $i\in U$ the coordinate integral contributes
$(1+\alpha_i)_{h_i-1}(1+v)^{-h_i-\alpha_i}$.  For $j\notin U$
it contributes
$((1+v)^{-\alpha_j}-1)/\alpha_j$.
After $x=(1+v)^{-1}$, the remaining integral is
\[
 \int_0^1(1-x)^{m-2}x^{A_U}
            \prod_{j\notin U}\frac{x^{\alpha_j}-1}{\alpha_j}\,dx.
\]
It equals
$(m-2)!\,(\prod_{j\notin U}\Delta_{\alpha_j})f_m(A_U)$ by the
ordinary beta integral.  This proves the formula first where the
coordinate integrations are absolute.  The face estimate and the
removable finite differences extend both sides holomorphically to
zero.  At zero, $f_m^{(\ell)}(0)=
(-1)^\ell\ell!\mathfrak h_\ell(m)/(m-1)!$, which recovers
\eqref{sfj:eq:monomial} independently.
\end{proof}

\section{A depth-two polylogarithm evaluation of every repeated-mass face}
\label{sfj:sec:polylog}

The general collision expansion has a stronger closed consequence when
all but one mass agree.  It holds at arbitrary nonnegative displacement,
beyond the disk used by the collision series.

\begin{lemma}[One-dimensional reduction of a repeated-mass face]
\label{sfj:lem:repeated}
Let $F\in C^2([B,B+\delta])$, where $\delta\ge0$.  For $k\ge2$,
\begin{equation}\label{sfj:eq:one-dimensional}
 \boxed{\displaystyle
 I_k(F;\underbrace{B,\ldots,B}_{k-1},B+\delta)
  =(-1)^{k-1}\int_0^1
   \frac{d}{dt}\left\{\frac{F(B+\delta t)-F(B)}t\right\}
                            [-\log(1-t)]^{k-1}\,dt.}
\end{equation}
The quotient is filled by its limiting value at zero; the integral is
ordinary and absolutely convergent.
\end{lemma}
\begin{proof}
After translating $B$ and scaling $\delta$, it suffices to check a
polynomial $F$.  Constants and affine terms vanish on both sides.
For $F(t)=t^m$, $m\ge2$, only the composition supported on the single
nonzero mass survives in \eqref{sfj:eq:monomial}; its value is
$(-1)^{k-1}(k-1)!\mathfrak h_{k-1}(m)$.  The right side gives
\[
 (-1)^{k-1}(m-1)\int_0^1
           t^{m-2}[-\log(1-t)]^{k-1}\,dt,
\]
which has exactly the same value by the beta coefficient calculation.
For general $C^2$ functions, approximate $F''$ uniformly by
polynomials and integrate twice, preserving $F(B)$ and $F'(B)$.
The face estimate gives continuity of the left side in the $C^1$
norm.  On the right,
\[
 \frac{d}{dt}\frac{F(B+\delta t)-F(B)}t
       =\delta^2\int_0^1 uF''(B+\delta tu)\,du,
\]
so uniform $C^2$ convergence and integrability of every fixed power of
$-\log(1-t)$ give continuity there as well.  This proves the lemma.
\end{proof}

We use the strict-index double polylogarithm
\begin{equation}\label{sfj:eq:double-polylog}
 \Li_{1,p}(z)=\sum_{m>n\ge1}\frac{z^m}{m n^p}
    =\sum_{m=2}^\infty\frac{H_{m-1}^{(p)}}m z^m,
     \qquad |z|<1.
\end{equation}
In particular
$\frac{d}{dz}\Li_{1,p}(z)=\Li_p(z)/(1-z)$.

\begin{theorem}[Every repeated-mass Gamma face has depth two]
\label{sfj:thm:gamma-polylog}
Let $a>0$, $B\ge0$, $\delta\ge0$, and $k\ge2$.  With $G_a$ as in
\eqref{sfj:eq:gammaG},
\begin{equation}\label{sfj:eq:gamma-polylog}
 \boxed{\displaystyle
 I_k(G_a;\underbrace{B,\ldots,B}_{k-1},B+\delta)
  =(-1)^{k-1}(k-1)!
     \sum_{n=0}^\infty
       \Li_{1,k-1}\!\left(
             \frac{\delta}{(n+a)^2+B+\delta}\right).}
\end{equation}
The sum is ordinary and absolutely convergent.  It converges normally
when $(a,B,\delta)$ ranges over a compact subset of
$\{a>0,B\ge0,\delta\ge0\}$.  Thus the dimension of the face does
not force polylogarithm depth greater than two in this parameter sector.
\end{theorem}

\begin{proof}
The Gamma product gives the normally convergent logarithmic difference
\begin{equation}\label{sfj:eq:gamma-difference}
 G_a(B+\delta t)-G_a(B)
       =-\sum_{n=0}^\infty
          \log\left(1+\frac{\delta t}{(n+a)^2+B}\right).
\end{equation}
For $u\ge0$, define
\[
 J_k(u)=\int_0^1
  \frac{\log(1+ut)-ut/(1+ut)}{t^2}
                            [-\log(1-t)]^{k-1}\,dt.
\]
Its integrand is nonnegative; the apparent singularity at $t=0$ is
removable.  Differentiation under the integral is justified on compact
$u$-intervals and gives
\[
 J_k'(u)=u\int_0^1\frac{[-\log(1-t)]^{k-1}}{(1+ut)^2}\,dt.
\]
Set $q=u/(1+u)$ and $y=1-t$.  Expanding $(1-qy)^{-2}$, with
nonnegative terms, yields
\[
 J_k'(u)
  =\frac{(k-1)!}{1+u}\Li_{k-1}\!\left(\frac{u}{1+u}\right).
\]
Since
\[
 \frac{d}{du}\Li_{1,k-1}\!\left(\frac{u}{1+u}\right)
      =\frac1{1+u}\Li_{k-1}\!\left(\frac{u}{1+u}\right)
\]
and both functions vanish at zero, we have
$J_k(u)=(k-1)!\Li_{1,k-1}(u/(1+u))$.
Insert \eqref{sfj:eq:gamma-difference} into
Lemma~\ref{sfj:lem:repeated}; positivity permits Tonelli's theorem,
and this scalar identity gives \eqref{sfj:eq:gamma-polylog}.

For the convergence assertion, all arguments lie in $[0,q_0]$, where
$q_0=\delta/(a^2+B+\delta)<1$.  From
$H_{m-1}^{(p)}/m\le1$ for $p\ge1$,
\[
 0\le\Li_{1,p}(q)\le\frac{q^2}{1-q}.
\]
The summands are therefore $O((n+a)^{-4})$, uniformly on the stated
compact parameter sets.  This proves finiteness as well as normal
convergence.  The case $\delta=0$ is immediate.
\end{proof}

For $k=2$ this is the previously known log-square covariance,
since $\Li_{1,1}(q)=\frac12\log^2(1-q)$.  For $k=3$, the double
polylogarithm reduces entirely to classical dilogarithms and
trilogarithms:
\begin{equation}\label{sfj:eq:Li12-classical}
 \Li_{1,2}(q)=2\Li_3(1-q)-2\zeta(3)
            -\log(1-q)\{\Li_2(1-q)+\zeta(2)\},\qquad 0<q<1.
\end{equation}
The formula follows by differentiating both sides and taking the
limit at $q=0$.  It provides an exact classical-polylogarithm lattice
evaluation of the first genuinely higher face.

For example, at $a=1$, $B=0$, and $\delta=3/10$,
\begin{equation}\label{sfj:eq:numeric-example}
 I_3(G_1;0,0,3/10)
   =2\sum_{n=1}^\infty\Li_{1,2}\!\left(\frac{3}{10n^2+3}\right)
   =0.0727899864111738847269555921654165064396\ldots.
\end{equation}
The decimal is a diagnostic value, not an interval enclosure.

\begin{remark}[A weighted extension]
The same proof applies when $G_a$ is replaced by the first spectral
jet $F_1(B)$ of $\sum_{n\ge0}W(n+a)((n+a)^2+B)^{-s}$ and
$\deg W\le2$.  Its second mass derivative is the ordinary series
$\sum W(n+a)((n+a)^2+B)^{-2}$.  More explicitly, with $x_n=n+a$,
\[
 F_1(B+\delta t)-F_1(B)-\delta tF_1'(B)
 =-\sum_{n=0}^{\infty}W(x_n)
   \left\{\log\left(1+\frac{\delta t}{x_n^2+B}\right)
                       -\frac{\delta t}{x_n^2+B}\right\}.
\]
This twice-anchored sum and its first two $t$ derivatives converge
normally on compact displacement intervals, since their summands are
$O(x_n^{\deg W-4})$.  Equality follows from the common second
derivative and the two zero initial values at $t=0$.  Every face
annihilates the two displayed affine terms.  The right side of
\eqref{sfj:eq:gamma-polylog}
then acquires the factor $W(n+a)$.  Its convergence follows from
$W(n+a)=O(n^2)$.  At higher polynomial degree, additional subtractions
are necessary and are not silently included in this statement.
\end{remark}

\section{Independent orders: the exact polar obstruction and mixed jets}
\label{sfj:sec:independent}

The common-ray family can be holomorphic at zero for every fixed ray
even when the joint family is not a holomorphic germ.  This is a
material distinction in extending the predecessor's theorem.

\subsection{A complete holomorphy criterion}

Fix $a>0$, $r\ge2$, masses $B_i\ge0$, and a polynomial
$W(x)=\sum_{d=0}^Dw_dx^d$.  Put $x_n=n+a$, $A=\sum_i\alpha_i$,
and initially, for $\Ree A>(D+1)/2$, set
\begin{equation}\label{sfj:eq:joint}
 \mathscr Z_W(\boldsymbol\alpha)
    =\sum_{n=0}^\infty W(x_n)
                  \prod_{i=1}^r(x_n^2+B_i)^{-\alpha_i}.
\end{equation}
All quadratic factors are positive, and their powers use real
logarithms.  Define the finite polynomials
\begin{equation}\label{sfj:eq:polar-poly}
 C_h(\boldsymbol\alpha)
  =[v^h]\prod_{i=1}^r(1+B_iv)^{-\alpha_i},\qquad
 P_W(\boldsymbol\alpha)=\sum_{h\ge1}w_{2h-1}C_h(\boldsymbol\alpha).
\end{equation}
Only $h\le(D+1)/2$ occurs in $P_W$.

\begin{theorem}[The complete joint polar part]
\label{sfj:thm:pole}
Near $\boldsymbol\alpha=0$ there is a holomorphic function
$\mathscr H_W$ such that
\begin{equation}\label{sfj:eq:pole}
 \boxed{\displaystyle
 \mathscr Z_W(\boldsymbol\alpha)
     =\frac{P_W(\boldsymbol\alpha)}{2A}
                         +\mathscr H_W(\boldsymbol\alpha),\qquad
 \mathscr H_W(0)=\sum_{d=0}^D w_d\zeta(-d,a).}
\end{equation}
The quotient is a prescribed polar subtraction; its part divisible
by $A$ is of course holomorphic.  The joint spectrum is holomorphic
at zero if and only if either $W$ is even or all the masses $B_i$
are equal.
\end{theorem}

\begin{proof}
Choose a finite head so that $B_i/x_n^2<\rho<1$ on the tail.  A
finite binomial expansion with an analytic remainder gives
\[
 \prod_i(x_n^2+B_i)^{-\alpha_i}
     =x_n^{-2A}\sum_{h=0}^{K}C_h(\boldsymbol\alpha)x_n^{-2h}
                      +\text{remainder}.
\]
Take $K$ sufficiently large that the weighted remainder and all
required parameter derivatives converge normally near zero; any
$K$ with $2K+2>D+1$ leaves a strict convergence margin after a
small neighborhood is chosen.  The removed terms are finite sums of
$w_dC_h(\boldsymbol\alpha)\zeta(2A+2h-d,a+N)$, plus the retained
head.  Only $2h-d=1$ gives a pole through the origin, and its
residue is exactly $w_dC_h/(2A)$.  Subtracting these terms proves
\eqref{sfj:eq:pole}.  At zero, $C_0=1$ and $C_h(0)=0$ for $h\ge1$;
the finite head and tail then give the displayed value of
$\mathscr H_W(0)$.

If $W$ is even, the polar polynomial vanishes.  If every mass is
$B$, then $C_h=(-1)^h(A)_hB^h/h!$, which is divisible by $A$ for
$h\ge1$.  These prove sufficiency.  Conversely, suppose $w_{2H-1}$
is the highest nonzero odd coefficient and $B_i\ne B_j$.  Restrict
to $\alpha_i=t$, $\alpha_j=-t$, and all other parameters zero.
Then $A=0$ and
\[
 \prod_\ell(1+B_\ell v)^{-\alpha_\ell}
  =\exp\{t[\log(1+B_jv)-\log(1+B_iv)]\}.
\]
The coefficient of $t^H$ in $C_H$ is
$(B_j-B_i)^H/H!$.  Lower $C_h$ have degree at most $h<H$ and
cannot cancel it.  Thus $P_W$ is not divisible by $A$, proving
necessity.
\end{proof}

For a fixed vector $\boldsymbol\lambda$ with
$L=\sum_i\lambda_i\ne0$, every straight ray has the finite limit
\begin{equation}\label{sfj:eq:direction-limit}
 \lim_{t\to0}\mathscr Z_W(t\boldsymbol\lambda)
 =\sum_dw_d\zeta(-d,a)+\frac1L\sum_i\lambda_i p_W(B_i),\qquad
 p_W(B)=\sum_{h\ge1}\frac{(-1)^hw_{2h-1}}{2h}B^h.
\end{equation}
This follows from the linear coefficient
$C_h(t\boldsymbol\lambda)=
t(-1)^h\sum_i\lambda_iB_i^h/h+O(t^2)$.
For example, $a=1$, $W=x$, and $(B_1,B_2)=(0,1)$ give
\begin{equation}\label{sfj:eq:direction-counterexample}
 \lim_{t\to0}\mathscr Z_x(t,t)=-\frac13,
 \qquad
 \lim_{t\to0}\mathscr Z_x(t,2t)=-\frac5{12}.
\end{equation}

Even agreement of all straight-ray limits does not suffice.  For
$W=x^3+x/2$ and $(B_1,B_2)=(0,1)$, the polar term is
$\alpha_2^2/(4(\alpha_1+\alpha_2))$.  Every non-tangential straight
ray has the same zero-order limit, but along
$\alpha_2=t$, $\alpha_1=-t+t^3$ this term is $1/(4t)$.
The joint germ is therefore singular.  These are obstructions to an
unqualified extension, not errors in the correctly scoped common-ray
theorem of the predecessor.

\begin{corollary}[Periodic mean cancellation]
If a bounded periodic weight $c_n$ is inserted in every sum, with
mean $\bar c$, the polar term in \eqref{sfj:eq:pole} is multiplied
by $\bar c$.  The joint germ is holomorphic exactly when
$\bar c=0$, or $W$ is even, or all masses agree.
\end{corollary}
\begin{proof}
For a period $q$, the auxiliary Dirichlet series is
$q^{-z}\sum_{v=0}^{q-1}c_v\zeta(z,(a+v)/q)$ and its residue at
$z=1$ is $\bar c$.  All other parts of the proof are unchanged.
\end{proof}

\subsection{The holomorphic mixed defect after an explicit subtraction}

Write
\[
 E_W(s;B)=\sum_{n\ge0}W(x_n)(x_n^2+B)^{-s},\qquad
 F_j(B)=\left.\partial_s^jE_W(s;B)\right|_{s=0}.
\]
The one-mass functions are holomorphic at zero and
$F_0(B)=\sum_dw_d\zeta(-d,a)+p_W(B)$.  Define
\begin{equation}\label{sfj:eq:joint-defect}
 \mathscr D_W(\boldsymbol\alpha)
   =\frac1A\sum_i\alpha_iE_W(A;B_i)
                           -\mathscr Z_W(\boldsymbol\alpha).
\end{equation}
Let $p_h=(-1)^hw_{2h-1}/(2h)$, and set
\begin{equation}\label{sfj:eq:joint-local}
 \mathscr L_W(\boldsymbol\alpha)
  =\sum_{h\ge1}p_h\left\{
    \frac{\sum_i\alpha_iB_i^h}{A}
    -\frac{h!}{(A)_h}
       \sum_{h_1+\cdots+h_r=h}
          \prod_i\frac{(\alpha_i)_{h_i}B_i^{h_i}}{h_i!}
                         \right\}.
\end{equation}
This is a finite rational function.  Its $h=1$ term is identically
zero.  In several variables its apparent origin singularities should
\emph{not} in general be filled by a common-ray value.  For instance,
\begin{equation}\label{sfj:eq:joint-local-cubic}
 \mathscr L_{x^3}(\boldsymbol\alpha)
   =\frac{\sum_{i<j}\alpha_i\alpha_j(B_i-B_j)^2}{4A(1+A)}.
\end{equation}

\begin{theorem}[Independent-order face decomposition]
\label{sfj:thm:joint-faces}
Put
\[
 \mathcal G_W(u;B)=\frac{E_W(u;B)-F_0(B)}u
       =\sum_{j\ge1}\frac{F_j(B)}{j!}u^{j-1},
 \qquad
 Q_J(\boldsymbol\alpha)=
 \frac{\Gamma(1+\sum_{i\in J}\alpha_i)}
                      {\prod_{i\in J}\Gamma(1+\alpha_i)}.
\]
The difference $\widetilde{\mathscr D}_W=
\mathscr D_W-\mathscr L_W$ is holomorphic near zero and equals
\begin{equation}\label{sfj:eq:joint-faces}
 \boxed{\displaystyle
 \widetilde{\mathscr D}_W(\boldsymbol\alpha)
   =-\sum_{\substack{J\subseteq\{1,\ldots,r\}\\|J|\ge2}}
       \left(\prod_{i\in J}\alpha_i\right)Q_J(\boldsymbol\alpha)
       \int_{\Sigma_J}\mathcal H_J(\mathcal G_W(A;\cdot);t)
                              \prod_{i\in J}t_i^{\alpha_i-1}\,dt.}
\end{equation}
A sufficient neighborhood for all displayed integrals and zero-point
germs is $\max_i|\alpha_i|<1/(4r)$.  The formula supplies ordinary,
absolutely convergent integrals and locally normal derivatives there.
\end{theorem}

\begin{proof}
For initially positive real parts, the usual Dirichlet parameterization
of the product in \eqref{sfj:eq:joint} has total order $A$.
Finite subset inversion decomposes the perspective into its faces.
Under the Dirichlet measure with parameters $\alpha_i$, the mass
$t_J$ has mean $A_J/A$, where $A_J=\sum_{i\in J}\alpha_i$, and
is independent of the normalized coordinates within $J$.  Thus the
contribution of face $J$ is
\[
 \frac{A_J}{A}\frac{\Gamma(A_J)}{\prod_{i\in J}\Gamma(\alpha_i)}
 \int_{\Sigma_J}\mathcal H_J(E_W(A;\cdot);t)
                                         \prod_{i\in J}t_i^{\alpha_i-1}\,dt
 =\frac{\prod_{i\in J}\alpha_i}{A}Q_J
                    \int_{\Sigma_J}\mathcal H_J(E_W(A;\cdot);t)
                                         \prod_{i\in J}t_i^{\alpha_i-1}\,dt.
\]
The singleton terms give the weighted endpoint average in
\eqref{sfj:eq:joint-defect}.  Splitting $E_W(A;B)=F_0(B)+
A\mathcal G_W(A;B)$ yields \eqref{sfj:eq:joint-faces}; ordinary
Dirichlet polynomial moments give exactly \eqref{sfj:eq:joint-local}
for the $F_0$ part.  The identity continues meromorphically from its
common absolute-convergence domain.  In the stated polydisk,
$|A|<1/4$, each one-mass zero germ is holomorphic, and
\eqref{sfj:eq:normal-face-domain} permits every face integral and
its derivatives.  This proves the claimed joint holomorphy.
\end{proof}

For even $W$, the local term is zero and both the joint spectrum and
the joint defect are holomorphic.  More generally, the \emph{raw
defect} $\mathscr D_W$ is holomorphic at zero precisely when the odd
part of $W$ has degree at most one, or all masses agree.  With periodic
weights, mean zero is an additional sufficient and necessary alternative.
Indeed its polar numerator is the negative of the degree-at-least-two
part of $P_W$; the highest-degree restriction used in
Theorem~\ref{sfj:thm:pole} proves necessity.  This criterion explains
why the spectrum and its defect have different analytic boundaries.

\subsection{A particularly simple formula for distinct mixed derivatives}

For a set $J$ put
\[
 \Pi_J(\boldsymbol\alpha;t)=
 \exp\left\{\sum_{i\in J}\alpha_i\log t_i+
  \sum_{n\ge2}\frac{(-1)^n\zeta(n)}n
    \left[\left(\sum_{i\in J}\alpha_i\right)^n-
                             \sum_{i\in J}\alpha_i^n\right]\right\}.
\]
This is $Q_J\prod_{i\in J}t_i^{\alpha_i}$.  If $\boldsymbol\nu$
is a nonnegative multiindex, extraction from
\eqref{sfj:eq:joint-faces} gives the finite identity
\begin{align}\label{sfj:eq:all-mixed}
 \partial^{\boldsymbol\nu}\widetilde{\mathscr D}_W(0)
 ={}&-\boldsymbol\nu!\!
 \sum_{\substack{J\subseteq\operatorname{supp}\boldsymbol\nu\\|J|\ge2}}
 \sum_{j=1}^{|\boldsymbol\nu|-|J|+1}\frac1{j!}
 \int_{\Sigma_J}\mathcal H_J(F_j;t)
       [\boldsymbol\alpha^{\boldsymbol\nu-\boldsymbol1_J}]
        \{A^{j-1}\Pi_J(\boldsymbol\alpha;t)\}
                          \frac{dt}{\prod_{i\in J}t_i}.
\end{align}
Thus a face contributes only if each of its labels is differentiated
at least once.  The inner jet has order at most
$|\boldsymbol\nu|-|J|+1$.

\begin{corollary}[Squarefree mixed derivatives]
\label{sfj:cor:squarefree}
If $U$ contains $m\ge2$ distinct labels, then
\begin{equation}\label{sfj:eq:squarefree}
 \boxed{\displaystyle
 \left(\prod_{i\in U}\partial_{\alpha_i}\right)
             \widetilde{\mathscr D}_W(0)
 =-\sum_{\substack{J\subseteq U\\|J|\ge2}}
           \frac{I_J(F_{m-|J|+1})}{m-|J|+1}.}
\end{equation}
No logarithmic simplex weights or Gamma Taylor corrections remain.
For three labels this is
\[
 \partial_{\alpha_1}\partial_{\alpha_2}\partial_{\alpha_3}
          \widetilde{\mathscr D}_W(0)
 =-\tfrac12\{I_{12}(F_2)+I_{13}(F_2)+I_{23}(F_2)\}
                                      -I_{123}(F_1).
\]
\end{corollary}
\begin{proof}
In \eqref{sfj:eq:all-mixed}, every variable belonging to $J$ is
already supplied by the face prefactor.  The remaining variables are
outside $J$, whereas $\Pi_J$ contains only variables in $J$.
Consequently it contributes only its constant coefficient.  The
required monomial is supplied by $A^{m-|J|}$ and has coefficient
$(m-|J|)!$.  Division by $(m-|J|+1)!$ gives the stated factor.
\end{proof}

For $\deg W\le2$, this has an ordinary logarithmic-sum interpretation.
Put $L_i(n)=\log((n+a)^2+B_i)$.  With $r\ge2$ labels,
\begin{equation}\label{sfj:eq:ordinary-log-identity}
 \sum_{n=0}^\infty W(n+a)
    \left\{\frac1r\sum_{i=1}^rL_i(n)^r-\prod_{i=1}^rL_i(n)\right\}
 =(-1)^{r+1}
      \sum_{\substack{J\subseteq\{1,\ldots,r\}\\|J|\ge2}}
           \frac{I_J(F_{r-|J|+1})}{r-|J|+1}.
\end{equation}
The leading and first-order changes in $B_i/(n+a)^2$ cancel in the
braces.  The summand is therefore
$O(n^{\deg W-4}\log^{r-2}(2+n))$, and the sum is absolutely
convergent.  This same estimate locally in the order parameters
justifies termwise differentiation of the defect.  The local term is
zero at these polynomial degrees.  Taking the squarefree derivative
in the pointwise defect proves \eqref{sfj:eq:ordinary-log-identity}.

\subsection{Verification and the remaining scope}

The supplied exact implementation compares
\eqref{sfj:eq:monomial} against a separately constructed rational
Dirichlet-moment formula for $2\le k\le5$ and $2\le m\le8$;
all 28 cases agree exactly over the rational numbers.  Six additional
quadratic variance checks cover $2\le k\le7$.

The numerical program uses 48 decimal digits.  It compares direct
triangle face integrals of degrees $3,4,6$ against the finite formula;
their absolute residuals are respectively approximately
$1.02\cdot10^{-26}$, $3.42\cdot10^{-25}$, and
$1.33\cdot10^{-23}$.  The repeated-mass one-dimensional formula and
the independently expanded positive double-polylogarithm series agree
for $k=2,\ldots,6$ to residuals below $1.7\cdot10^{-38}$.
Most significantly, the original \emph{two-dimensional Gamma face}
in \eqref{sfj:eq:numeric-example} was evaluated directly on all six
coordinate-ordering sectors.  At Gauss--Legendre orders $16,24,32$,
its residual against the polylogarithm lattice value decreases from
$2.33\cdot10^{-19}$ to $3.00\cdot10^{-28}$ and then
$4.41\cdot10^{-37}$.  The sector map is
$t=(1,u,uv)/(1+u+uv)$ and its weighted measure is $du\,dv/(uv)$.
The face numerator is retained as a combined expression throughout.

These computations are finite algebraic checks and numerical
diagnostics, not interval certificates or substitutes for the proofs.
The results complete the specified all-order collision and
independent-order extensions relative to the inspected spectral-face
report, and give a global depth-two evaluation in the one-displaced-mass
sector.  They do not prove a finite classical-polylogarithm reduction
for arbitrary distinct masses or every polylogarithm weight.  Natural
next targets are a finite dictionary for two distinct mass clusters,
arithmetic reductions of the resulting double-polylogarithm lattice
sums, and a basis-independent description of the joint polar
subtraction when the affine spectral factors have different degrees.
